\section{Related Works}
\label{RelatedWorks}
\textbf{Efficient Visual Tracking.}
Visual object tracking aims to localize a target object in each video frame based on its initial appearance. Tracking methods have evolved from hand-crafted features~\cite{KCF, DCFST, LADCF} to deep learning-based approaches~\cite{SiamFC, SiamRPN, Ocean, ATOM, SiamRPN++}. Earlier models adopt a two-stream pipeline that decouples feature extraction and relation modeling. Recently, Transformers~\cite{vit} have revolutionized visual tracking by enabling more effective feature interaction. The emergence of one-stream pipelines~\cite{ostrack, mixformer, mixformerv2, artrack, simtrack} integrates feature extraction and relation modeling into a unified process, significantly improving tracking performance. However, the quadratic complexity of ViT self-attention introduces a significant bottleneck for real-time inference, especially on resource-constrained devices.
To address this efficiency bottleneck, researchers have explored several strategies. Early works~\cite{lighttrack} employed neural architecture search to design lightweight Siamese networks, but the high cost of the search process limits their practicality. Other approaches~\cite{fear, HCAT, ETTrack, HiT, LiteTrack, AVTrack} focus on designing compact models by reducing the number of parameters. However, these methods often suffer from limited model capacity, restricting their accuracy on challenging benchmarks. Model compression techniques such as layer pruning~\cite{LiteTrack, mixformerv2} and knowledge distillation~\cite{compressTracker, mixformerv2} have also been explored, but pruning often causes significant accuracy degradation. Furthermore, some trackers~\cite{AVTrack, wu2024learning, abavitrack} introduced conditional dynamic computation to selectively process tokens or layers based on predefined thresholds, achieving better \emph{average} running times. However, the resulting variable per-frame computation complicates hardware scheduling and does not reduce the worst-case inference latency. In contrast, our approach uses fixed top-$K$ token selection, ensuring predictable and consistent latency across all frames.

\textbf{Temporal Priors in Visual Tracking.}
Temporal priors, especially motion locality, are well-established in visual tracking: target motion between consecutive frames is typically small relative to the search region. Classical approaches exploit this prior via a penalty window (Hanning/cosine) applied to the output score map. SiamFC~\cite{SiamFC} first applied a cosine window to the cross-correlation response to suppress high-response regions near the boundary of the search area. SiamRPN++~\cite{SiamRPN++} further combines a cosine window penalty with a size-change penalty on the regression response to reduce large, abrupt displacement predictions. Discriminative correlation filters such as KCF~\cite{KCF} and ECO~\cite{ECO} apply a cosine window to the input features to reduce periodic boundary artefacts, which also implicitly penalises predictions near the edges. Many modern trackers, including DiMP~\cite{DiMP}, prDiMP~\cite{prDIMP}, STARK~\cite{STARK}, and OSTrack~\cite{ostrack}, likewise multiply a Hanning window onto the predicted score map as a post-processing step before selecting the final target location.
Beyond hand-crafted windows, recent Transformer-based trackers also explore data-driven temporal modeling by encoding historical predictions as tokens and letting the network learn to use this information, e.g., SwinTrack~\cite{lin2022swintrack} and the ARTrack series~\cite{artrack,wang2026fartrack}. However, these designs do not explicitly impose a spatial locality prior during computation.
Despite their widespread adoption, penalty windows are applied only after the full forward pass and merely re-weight already-computed scores, leaving the heavy encoder computation unaffected. In contrast, our approach injects temporal motion priors into the token reduction process itself, using them to guide which tokens are computed in subsequent layers and thereby reducing computation.

\textbf{Token Sparsification in Vision Transformers.}
The tokenized processing paradigm of Vision Transformers enables novel sparsification strategies that are absent in CNN-based trackers. EViT~\cite{evit} pioneered token pruning by selecting tokens with the highest attention to the class token. OSTrack~\cite{ostrack} adapted this idea by measuring cross-attention between search and template tokens, retaining the top-K most correlated tokens. However, their reliance on early-layer attention scores introduces noise. 
As a result, many methods apply token pruning conservatively and progressively at intermediate layers, leading to limited speedup.
DynamicViT~\cite{dynamicvit} proposed learnable importance predictors using lightweight MLPs, and DToP~\cite{DTP4SS} further explored this idea in semantic segmentation. Similarly, GRM~\cite{grm} applied such techniques to guide token interactions in tracking tasks. Although these methods improve robustness. AVTrack~\cite{AVTrack} also uses sub-network for early token exit, but the resulting variable computation per frame complicates hardware scheduling and deployment. In contrast, our approach combines motion-guided locality priors with fixed token budgets (top-K selection), ensuring predictable and consistent latency.